% Generated draft: requires mathematical and source review.
\documentclass[11pt]{article}
\usepackage[margin=25mm]{geometry}
\usepackage{fontspec}
\setmainfont{Latin Modern Roman}
\usepackage{amsmath,amssymb,mathtools,unicode-math}
\setmathfont{Latin Modern Math}
\usepackage{enumitem,needspace,xurl,hyperref}
\setlength{\parindent}{0pt}
\setlength{\parskip}{5pt}
\setlength{\emergencystretch}{3em}
\newcommand{\N}{\mathbb{N}}
\newcommand{\Z}{\mathbb{Z}}
\newcommand{\Q}{\mathbb{Q}}
\newcommand{\R}{\mathbb{R}}
\newcommand{\C}{\mathbb{C}}
\newcommand{\F}{\mathbb{F}}
\newcommand{\A}{\mathbb{A}}
\newcommand{\PP}{\mathbb{P}}
\newcommand{\E}{\mathbb{E}}
\newcommand{\eps}{\varepsilon}
\newcommand{\dd}{\,\mathrm{d}}
\begin{document}
\section*{UnsolvedMath ID: 701424}
\textbf{Source problem code:} catalogue entry
\subsection{Definitions and mathematical statement}
For a finite group \(G\), let its element-order spectrum be
\[
\pi_e(G)=\{\operatorname{ord}(g):g\in G\}.
\]
A composition factor of \(G\) is a simple quotient occurring in a composition series of \(G\). A section of a group \(K\) is a quotient \(H/N\), where \(H\le K\) and \(N\trianglelefteq H\).

The supplied statement refers to Archive problem 13.63 for definitions, but that entry is not included. In particular, the meaning of \(h\) cannot be confirmed from the supplied record. To make a conditional formulation explicit, for a finite set \(S\subseteq\Z_{>0}\), define
\[
h_{\mathrm{fin}}(S)=\#\{[K]:K\text{ is a finite group and }\pi_e(K)=S\},
\]
where \([K]\) denotes the isomorphism class of \(K\).

\textbf{Problem (conditional on the intended definition of \(h\)).} Under the interpretation \(h=h_{\mathrm{fin}}\), do there exist finite groups \(G_1,G_2\) such that
\[
\pi_e(G_1)=\pi_e(G_2),\qquad h_{\mathrm{fin}}(\pi_e(G_1))<\infty,
\]
and such that, for every \(i,j\in\{1,2\}\) with \(i\ne j\), every nonabelian composition factor \(F\) of \(G_i\), every subgroup \(H\le G_j\), and every normal subgroup \(N\trianglelefteq H\), one has
\[
F\not\cong H/N?
\]
The identification of the source's \(h\) with \(h_{\mathrm{fin}}\) requires verification against Archive 13.63. The supplied statement does not separately require \(G_1\not\cong G_2\) or require either group to have a nonabelian composition factor; no such hypotheses have been added.
\subsection{Short English statement}
Conditionally interpreting \(h\) as the number of finite-group isomorphism classes with a prescribed element-order spectrum, ask whether two finite groups can share a spectrum realized by only finitely many such classes while every nonabelian composition factor of either group fails to occur as a section of the other. The referenced definition of \(h\) remains to be checked.
\subsection{Sources}
\begin{itemize}
\item Local UnsolvedMath registry, permanent ID \(701424\), exact source problem code \texttt{catalogue entry}, title ``Kourovka Notebook Problem 17.123''; supplied upstream revision \texttt{d371ee917ac9647fbb06cd6fc529deec6f3dda82}.
\item The Kourovka Notebook (Unsolved Problems in Group Theory), problem 17.123, 17th issue (2010); proposed by W. J. Shi. The supplied background also identifies the 21st issue edition (2026), \texttt{arXiv:1401.0300v48}. The collection description names E. I. Khukhro and V. D. Mazurov as editors.
\item The source explicitly refers to Archive problem 13.63 for definitions. Its contents were not supplied, so those definitions have not been verified.
\item The registry supplies the status ``open'' and an update timestamp of \texttt{2026-10-16T14:00:00Z}. These are source metadata, not an independently verified literature assessment.
\end{itemize}
\end{document}
